\subsection{複数の meridian blow-up：homotopy 型と整数交点格子}
\label{cs7:meridian-section}

一つの node を取り込む場合の比較を、同じ成分上の複数の blow-up にも使える形にする。
この節の計算は実際の blow-up complement に対するものであり、任意の focus--focus
fiber の取り込みをこの模型で表せるという主張ではない。

$m\geq4$ 本の一般位置の直線 $L_1,\ldots,L_m\subset\PP^2$ を固定する。
$L_i$ の double points を避けた相異なる $r_i\geq0$ 点を blow up し、
得られる曲面を $M_{\boldsymbol r}$、直線の strict transforms の union を
$D_{\boldsymbol r}$、その補集合を $U_{\boldsymbol r}$ とする。
exceptional divisors は $D_{\boldsymbol r}$ に含めない。
\[
 r=\sum_i r_i,\qquad s=\#\{i:r_i>0\},\qquad
 n=m-1,\qquad k=m-1-s
\]
と書き、以下では $s<m$、すなわち少なくとも一本の直線を blow up しないと仮定する。
$T^k_{(2)}$ は標準積 CW 構造の $2$-skeleton を表し、$k\leq2$ では $T^k$ 自身とする。

\begin{proposition}[meridian attachments の homotopy 型]
\label{cs7:homotopy}
上の仮定の下で、
\begin{equation}\label{cs7:wedge}
 U_{\boldsymbol r}\simeq
 T^k_{(2)}\vee\bigvee^{\,t}S^2,
 \qquad
 t=\binom{m-1}{2}-\binom{k}{2}+r-s.
\end{equation}
従って
\[
 \pi_1(U_{\boldsymbol r})=\ZZ^k,\qquad
 b_2(U_{\boldsymbol r})=\binom{m-1}{2}+r-s.
\]
特に命題~\ref{cs6:invariants} の $U$ では
\begin{equation}\label{cs7:sextic-wedge}
 U\simeq T^3_{(2)}\vee S^2\vee S^2\vee S^2.
\end{equation}
$R=\ZZ[x_1^{\pm1},x_2^{\pm1},x_3^{\pm1}]$ とすると、その deck action を含めて
\begin{equation}\label{cs7:sextic-pi2}
 \pi_2(U)\cong R^4
\end{equation}
である。
\end{proposition}

\begin{proof}
各 exceptional divisor の $U_{\boldsymbol r}$ 内の部分は $\mathbb A^1$ であり、
これらは互いに交わらない。命題~\ref{cs6:invariants} と同じ tubular-neighborhood
pushout を各部分に適用すると、$U_{\boldsymbol r}$ は $W_m$ に $r_i$ 個ずつ
$L_i$ の meridian を attaching loop とする $2$-cells を付けた空間と homotopy 同値である。
これは相異なる点での blow-up であることを使い、infinitely near points は含めない。

blow up していない直線を infinity に選ぶと、残る $m-1$ 本の affine meridians は
Hattori の $T^{m-1}_{(2)}$ 模型の標準 generators に対応する。
ここでは単に cohomology ring を用いるのでなく、arrangement complement から
$(\CC^*)^{m-1}$ への defining-linear-functions map と整合する skeleton 模型を用いる。
この標識は Hattori~\cite[\S2, equation~(2.5), Proposition~(2.16)(2), pp.~207--213;
\S4, Lemma~(4.2), pp.~217--219]{Hattori} の構成に含まれる。
特に Proposition~(2.16)(2) は torus への inclusion との homotopy 可換性を保証し、
\S4 の local trivialization は各 labeled coordinate hyperplane を保存する。
一般位置 arrangement の skeleton 型の解説と meridian の cohomological marking については
Papadima--Suciu~\cite[\S\S1.1, 4.3; proof of Theorem~4.8(4)]{PapadimaSuciu}
も参照されたい。

この標準模型は一つの $0$-cell、$n$ 個の $1$-cells、および
$[x_i,x_j]$ を attaching words とする $\binom n2$ 個の $2$-cells を持つ。
殺す $s$ 個の generator について、それぞれ一本の circle と新しい一枚の disk を
合わせた contractible subcomplex を collapse する。
残る $k$ 個の generators の間の commutator cells は $T^k_{(2)}$ を与える。
殺した generator を含む commutator の attaching loop は null-homotopic となり、
それぞれ一つの $S^2$ を与える。同じ meridian に沿う二枚目以降の新しい cells も
同様に $S^2$ を与える。この数は $\binom n2-\binom k2+r-s$ である。
collapse した subcomplex は contractible な CW subcomplex なので、この操作は
homotopy 同値を与える。

$k=3$ のとき、$T^3$ の universal cover の cellular complex は
$x_1-1,x_2-1,x_3-1$ の Koszul complex である。
その最後の differential
\[
 d_3:R\longrightarrow R^3,\qquad
 1\longmapsto (x_3-1)e_{12}-(x_2-1)e_{13}+(x_1-1)e_{23}
\]
は単射であり、$\ker d_2=\operatorname{im}d_3\cong R$ である。
universal cover で追加の各 $S^2$ は自由 $R$ summand を一つずつ与える。
$m=5,r=s=1$ では $t=3$ なので \eqref{cs7:sextic-pi2} が従う。
\end{proof}

\begin{proposition}[整数交点格子と境界の torsion]
\label{cs7:lattice}
$K_{\boldsymbol r}\subset U_{\boldsymbol r}$ を divisor neighborhood の補集合として得られる
compact core、$Y_{\boldsymbol r}=\partial K_{\boldsymbol r}$ とする。
$H_2(K_{\boldsymbol r};\ZZ)$ の絶対 intersection lattice は
\begin{equation}\label{cs7:intersection}
 0^{\binom{m-1}{2}}\ \oplus\!
 \bigoplus_{i:r_i>0}(-A_{r_i-1}),
\end{equation}
である。$A_0=0$ とし、$A_q$ は対角 $2$、隣接成分 $-1$ の標準 root lattice とする。
特に
\[
 \sigma(K_{\boldsymbol r})=-(r-s).
\]
境界は plumbing matrix $J_m-\operatorname{diag}(r_1,\ldots,r_m)$ を持ち、
\begin{equation}\label{cs7:boundary-h1}
 H_1(Y_{\boldsymbol r};\ZZ)\cong
 \ZZ^{\binom{m-1}{2}+m-1-s}
 \oplus\bigoplus_{i:r_i>0}\ZZ/r_i\ZZ.
\end{equation}
その reduced JSJ graph は $K_m$ である。
境界包含は $H_1(Y_{\boldsymbol r})\to H_1(K_{\boldsymbol r})=\ZZ^k$ の全射を誘導し、
その kernel は
\[
 \ZZ^{\binom{m-1}{2}}\oplus\bigoplus_{i:r_i>0}\ZZ/r_i\ZZ
\]
に同型である。この抽象群の同型だけでは、境界内の具体的な peripheral subgroup は定まらない。
\end{proposition}

\begin{proof}
exceptional classes を $E_{i,a}$ と書けば
\[
 [D_i]=H-\sum_{a=1}^{r_i}E_{i,a}.
\]
pair $(M_{\boldsymbol r},U_{\boldsymbol r})$ の exact sequence と divisor に対する
duality により
\[
 \operatorname{im}\bigl(H_2(U_{\boldsymbol r})\to H_2(M_{\boldsymbol r})\bigr)
 =\ker\!\left(H_2(M_{\boldsymbol r})\xrightarrow{(\,\cdot D_i\,)_{i=1}^m}\ZZ^m\right).
\]
ここで $H_2(M_{\boldsymbol r},U_{\boldsymbol r})\cong H^2(D_{\boldsymbol r})\cong\ZZ^m$ である。
blow up していない成分との交点が零なので、kernel の $H$ coefficient は零である。
残る条件は各 $i$ について $\sum_a b_{i,a}=0$ となる。
従って kernel は
\[
 \bigoplus_{i:r_i>0}
 \left\{\sum_a b_{i,a}E_{i,a}:\sum_a b_{i,a}=0\right\},
\]
すなわち $\bigoplus_i(-A_{r_i-1})$ である。
この自由群への全射は整数加群として split し、補集合内の交点数は compactification
内の交点数と一致する。命題~\ref{cs7:homotopy} の $b_2$ から radical の rank は
$\binom{m-1}{2}$ となり、\eqref{cs7:intersection} を得る。

plumbing graph の各 edge は以前と同じ double point に対応し、各 diagonal weight は
$1-r_i$ である。よって matrix は表示されたものになる。
blow up していない一つの row を選ぶと、その relation は $\sum_i\mu_i=0$ である。
他の rows からこの row を引くと、残る relations は $r_i\mu_i=0$ となる。
未変更成分の meridian を一つ消去すれば
\[
 \operatorname{coker}Q\cong\ZZ^{m-1-s}\oplus\bigoplus_{r_i>0}\ZZ/r_i\ZZ.
\]
plumbing の exact sequence は、これに
$H_1(K_m;\ZZ)\cong\ZZ^{\binom{m-1}{2}}$ を加える。
後者は自由群なので sequence は群として split する。

各 Seifert piece の base は $(m-1)$ 穴 sphere で、その Euler 数は $3-m<0$ である。
隣接する normal-circle slopes の交点数は $1$ のままである。
従って命題~\ref{cs6:boundary} と同じ incompressibility と slope の議論により、
全ての plumbing tori が reduced JSJ system をなす。
最後の kernel の記述は、pair $(K_{\boldsymbol r},Y_{\boldsymbol r})$ の exact sequence
と $H^3(K_{\boldsymbol r})=0$ による全射、およびその自由群 $\ZZ^k$ への splitting から従う。
\end{proof}

\begin{remark}[Stein model の contact 不変量]
\label{cs7:contact}
この補集合は affine で、canonical line bundle は自明である。
実際、$N\gg0$ に対して $NH-\sum E_{i,\ell}$ は ample であり、
blow up した成分に係数 $1$、残る成分に正の整数係数を配分すれば、
この class は $D_{\boldsymbol r}$ を support とする ample divisor で表される。
また、未変更成分 $D_j=H$ を一つ選ぶと
\[
 K_{M_{\boldsymbol r}}=-\sum_{i:r_i>0}D_i+(s-3)D_j
\]
である。対応する Stein contact boundary には
\[
 c_1(\xi)=0,\qquad
 d_3(\xi)=\frac{-2\bigl(\binom{m-2}{2}+r\bigr)+3(r-s)}4
\]
が成り立つ\cite[Definition~4.2, Theorem~4.5]{Gompf}。
これも contactomorphism の十分条件ではない。
\end{remark}

\begin{corollary}[canonical pants 数に合わせる meridian 模型]
\label{cs7:canonical-budget}
$m\geq5$ とし、$P=(m-3)^2$ 個の標準四次元 pants と同じ Euler 数を持つように
$r=\binom{m-3}{2}$ 個の上記 blow-ups を行う。
さらに $\pi_1(U_{\boldsymbol r})=\ZZ^3$ を要求すると
\[
 s=m-4,\qquad
 U_{\boldsymbol r}\simeq T^3_{(2)}\vee\bigvee^{P-1}S^2,
 \qquad\pi_2(U_{\boldsymbol r})\cong\ZZ[\ZZ^3]^P,
\]
\begin{equation}\label{cs7:signature-budget}
 \sigma(K_{\boldsymbol r})=-\binom{m-4}{2}.
\end{equation}
$m\geq6$ なら、この compact core は disjoint full-face tree assembly の
$P$ 個の pants とは diffeomorphic でない。
\end{corollary}

\begin{proof}
命題~\ref{cs7:homotopy} より $m-1-s=3$ であり、同命題に $r,s$ を代入すると
wedge 数は $P-1$ となる。signature は
$r-s=\binom{m-4}{2}$ から従う。

full-face tree assembly について、face の $H_1$ からそれぞれの pants の $H_1$ への
写像は同型である。従って Mayer--Vietoris sequence を tree に沿って反復すると、
assembly の $H_2$ は各 pants に含まれる $H_2$ classes で生成される。
標準 pants の絶対交点形式は零であり、異なる pants に含まれる cycles は各 piece の
内部へ移せば互いに交わらない。よって assembly の絶対交点形式も零である。
$m\geq6$ では \eqref{cs7:signature-budget} は負なので矛盾する。
\end{proof}

例えば $m=6$、$(r_1,r_2)=(2,1)$ では $P=9$、
$U\simeq T^3_{(2)}\vee8S^2$ であるが、交点格子には $-A_1$ が現れ、
$H_1(Y)=\ZZ^{13}\oplus\ZZ/2\ZZ$ となる。
同じ homotopy 型の候補を比較しても、境界と交点格子は別に検査する必要がある。
これらの位相計算だけから Weinstein equivalence や exact symplectic surgery は従わない。
